\clearpage
\fancyhead[L]{\small KLÜBER STERNWARTE ZU MANNHEIM.}
\hypertarget{native-541-lower}{}\label{native-541-lower}
\begin{center}\small Göttingische gelehrte Anzeigen. Stück 192 u. 193. Seite 1913..1916. 1811 December 2.\\[3pt]\rule{.72\linewidth}{.3pt}\end{center}

\emph{Die Sternwarte zu Mannheim, beschrieben von ihrem Curator, dem Staats- und Kabinetsrath} \textsc{Klüber}. Mit einer Abbildung der Sternwarte in Steindruck. 1811. klein Folio 62 Seiten. In Commission bei \textsc{Gottlieb Braun}. Mannheim und Heidelberg.

{\rightskip=0pt plus 1em
Die Mannheimer Sternwarte ist gegenwärtig unstreitig von allen in Deutschland am reichsten mit Instrumenten versehen; Jeder, dem die Astronomie werth ist, nimmt ein lebhaftes, warmes Interesse an ihr, und es kann daher eine ausführliche Erzählung von ihrer Entstehung, ihrem Fortgange und ihrem gegenwärtigen Zustande nicht anders als willkommen sein. Doppelt erfreulich ist der Bericht aus der Feder des Curators der Sternwarte, dem der Ruhm und das Beste des Instituts so sehr am Herzen liegt, und der davon bereits so thätige Beweise gegeben hat.\par}

Der Grund zu dem durchaus massiven Gebäude wurde im Jahre 1772 unter dem Churfürsten von der Pfalz, \textsc{Carl Theodor}, nach dem Plane des bekannten Astonomen P. \textsc{Mayer} gelegt, und, den Bau vollkommen dauerhaft und schön auszuführen, wurden keine Kosten gespart; sie beliefen sich über 70000 Gulden. Man wird sich darüber nicht wundern, wenn man die ansehnliche Höhe des Gebäudes, 111 Fuss, erwägt, die damals für nothwendig gehalten wurden, aber freilich bei der gegenwärtigen Beobachtungsart als ein wesentlicher Fehler in der Anlage betrachtet werden muss, da sie das Beobachten eben so sehr erschwert, als der Zuverlässigkeit desselben schadet. Die Sternwarte wurde mit Instrumenten von den besten Engl. Künstlern versehen. Ein vortrefflicher achtfussiger Mauerquadrant von \textsc{Bird} wurde schon 1775 aufgestellt; im Jahre 1778 erhielt die Sternwarte einen zwölffussigen Zenith-Sector von \textsc{Sisson} und einen \textsc{Arnold}'schen Regulator. Späterhin lieferte \textsc{Ramsden} (für $145\tfrac{1}{2}$ Guineen) das sechsfussige Mittagsfernrohr. Ein dreifussiger Multiplicationskreis von \textsc{Reichenbach} in München wurde erwartet, und ist gegenwärtig, Privatnachrichten zufolge, bereits wirklich angekommen. Der erste Astronom, \textsc{Christian Mayer}, starb 1783 (sein durch Specialtafeln für Aberration und Nutation be\SourceLeaf{542}kannter Gehülfe, \textsc{Joh. Metzger}, starb 1780); vom Jahre 1784 bis 1786 stand der Sternwarte \textsc{Carl König}, in den Jahren 1786 und 1787 \textsc{Joh. Nepomuk Fischer} vor, dessen designirter Nachfolger, \textsc{Peter Ungeschick}, vor Antritt der Stelle starb; endlich seit 1788 ist \textsc{Roger Barry} (eine Zeit lang gemeinschaftlich mit \textsc{Henry}) Astronom der Sternwarte. Der Krieg war den Arbeiten sehr ungünstig, und zu Ende des Jahrs 1794 mussten alle Instrumente abgenommen und in Kisten eingepackt werden, worin sie 6 Jahre blieben Mehrere Male wurde die Sternwarte sogar bei dem Bombardement Mannheims von Haubitzen getroffen, ohne doch wesentlich beschädigt zu werden, und der Astronom sah sich persönlichen Misshandlungen und Verfolgungen ausgesetzt, denen er sich nur durch die Flucht entziehen konnte. Nach dem Lüneviller Frieden, wo Mannheim an das Haus Baden fiel, wurden zwar mehrere wissenschaftliche und Kunstsammlungen von Mannheim nach München gebracht, allein die kostbaren astronomischen Instrumente wurden jener Stadt erhalten und 1801 wieder aufgestellt. Erst später wurden die nothwendigsten Kosten zur Bestreitung der Unterhaltung der Sternwarte und zur Besoldung des Astronomen angewiesen.

Im Jahre 1808 wurde die Sternwarte der nähern Fürsorge des Staatsraths \textsc{Klüber} übergeben, der alles that, ihr wieder aufzuhelfen. Die erforderlichen Kosten zur Unterhaltung, Verschönerung und Möblirung des Gebäudes, um auch gelegentlich fremden Astronomen zum Aufenthalt dienen zu können, ferner für Beleuchtung, Heizung, für Correspondenz, für eine Bibliothek und andere kleine Bedürfnisse wurden bewilligt: ein schöner Obelisk zur Mire méridienne wurde auf der Nordseite aufgeführt, und ein zweiter wird auf der Südseite gegenwärtig errichtet. Zum Druck des grossen Stern-Catalogs, an welchem \textsc{Barry} seit vielen Jahren arbeitet, hat die Regierung die Kosten angewiesen, ein Multiplicationskreis von \textsc{Reichenbach} ist, wie bereits oben bemerkt wurde, erst neulich angeschafft, und dem Astronomen selbst wurde eine ansehnliche Gehaltserhöhung zu Theil, und ihm zugleich ein besoldeter Aufwärter beigegeben, so wie man jetzt auf einen Adjuncten für den Astronomen bedacht ist.

Gewiss verdient die Badensche Regierung und der würdige Curator der Sternwarte wegen dieser freigebigen Unterstützung einer schönen, der praktischen Astronomie gewidmeten, Anstalt den Dank aller Freunde dieser Wissenschaft, und man darf mit Recht von einer thätigen und einsichtsvollen Benutzung derselben reiche Früchte hoffen.

Der übrige Theil der vorliegenden Schrift gibt noch eine Uebersicht über die bisherigen Ortsbestimmungen der Sternwarte, deren Länge von Paris auf $24^{\mathrm m}31^{\mathrm s}8$, so wie die Breite auf $49^\circ\,29'\,16''$ festgesetzt wird; ein Verzeichniss der sämmtlichen vorhandenen Instrumente; biographische Nachrichten von den bisherigen Astronomen der Sternwarte, und ein vollständiges Verzeichniss ihrer Schriften.
